\section{Experiment 1: Health Insurance}

Focus on \citet{johnson2013can}.

The experiment manipulates the number of options shown to the subjects: 4 versus 8.

We need to focus on the following metrics:
Choice accuracy: modelled using logistic regression - Normatively optimal choice.
Dollar error (regret): modelled using ANOVA

Bias measurement:
Some attributes overweigh others, such as policy premiums versus deductibles.

Confidence rating (1-9 scale): 
compare with the actual performance see what the correlation is.
Low correlation between confidence and actual performance shows overconfidence bias.

Decision, time, and effort measure through number of tokens generated and amount of budget consumed.

Optional: The effect of extra tools such as calculator on the decision quality and confidence.
Optional: Asking the model to verify its choice and see if the confidence increases? If that's the case, does the increase in confidence correlate with the final performance?
In other words, can we ask the model to evaluate its own performance (LLM as a judge)?


We need to compare choice accuracy and dollar error to the human benchmarks from the paper.
compare LLM confidence with human benchmarks 


Managerial Implications: This experiment helps policymakers understand how many billions of dollars saved/lost annually based on the errors that are made?